% GENERATED FILE. Edit canonical lessons, then run npm run generate:books.
% canonical-source-hash: fnv1a64:f1ec24ff9bc8cc12
% canonical-lessons: UR-C343-mahkama, UR-C343-darkhast, UR-C343-uhda, UR-C343-hidayat, UR-C343-sujhana

\chapter{Offices and Officials}
\label{ch:ur-offices-and-officials}
\hlchaptermodality{343}

\begin{chapteropening}
\textbf{By the end of this chapter:} \emph{I can deal with an office and its officials.}

Complete the last lesson of chapter 343: I can deal with an office and its officials.
\end{chapteropening}

\section[\texorpdfstring{\ur{محکمہ} (mahkamā)}{محکمہ (mahkamā)}]{\ur{محکمہ} (mahkamā) — a government department}

\label{lesson:UR-C343-mahkama}



\begin{quote}
Before the new one: say the Urdu for a candidate, then the Urdu for to impose.
\end{quote}

\subsection*{You'll want to know: \ur{محکمہ}}
\textbf{\ur{محکمہ}} — \emph{mahkamā} — \textquotedblleft{}a government department\textquotedblright{}.

\begin{quote}
\ur{وہ} \ur{پولیس} \ur{کے} \ur{محکمے} \ur{میں} \ur{کام} \ur{کرتا} \ur{ہے۔}
\end{quote}

\begin{quote}
\emph{voh polīs ke mahkame meṅ kām kartā hai.}
\end{quote}

\begin{quote}
He works in the police department.
\end{quote}

Masculine.

\begin{cousinweb}
\textbf{\ur{محکمہ}} (\emph{mahkamā}) is usually traced, through Persian, to Arabic \emph{maḥkama}, \textquotedblleft{}a court\textquotedblright{}; Urdu uses it for a department of government, a sense our sources do not trace.
\end{cousinweb}

\begin{grammarlens}[title={What you've built}]
One more word for giving reasons and opinions.
\end{grammarlens}

\subsection*{Your turn}
\begin{itemize}
\raggedright
  \item \emph{Say it:} \emph{mahkamā}
  \item \emph{Say it:} \emph{mahkamā}, once more
  \item \emph{Say it:} \emph{voh polīs ke mahkame meṅ kām kartā hai.} — He works in the police department.
  \item \emph{Recall:} say the Urdu for a candidate, then the Urdu for to impose, then say \emph{mahkamā} again
\end{itemize}

\begin{tcolorbox}[breakable,colback=teal!4,colframe=teal!35!black,title={Before you move on}]
What does \ur{محکمہ} mean? (\textquotedblleft{}A government department\textquotedblright{}.) Say it once more.
\end{tcolorbox}
\section[\texorpdfstring{\ur{درخواست} (darkhāst)}{درخواست (darkhāst)}]{\ur{درخواست} (darkhāst) — an application; a request}

\label{lesson:UR-C343-darkhast}



\begin{quote}
Before the new one: say the Urdu for to impose, then the Urdu for a government department.
\end{quote}

\subsection*{You'll want to know: \ur{درخواست}}
\textbf{\ur{درخواست}} — \emph{darkhāst} — \textquotedblleft{}an application; a request\textquotedblright{}.

\begin{quote}
\ur{میں} \ur{نے} \ur{نوکری} \ur{کے} \ur{لیے} \ur{درخواست} \ur{دی۔}
\end{quote}

\begin{quote}
\emph{maiṅ ne naukrī ke liye darkhāst dī.}
\end{quote}

\begin{quote}
I applied for a job.
\end{quote}

Feminine.

\begin{cousinweb}
\textbf{\ur{درخواست}} (\emph{darkhāst}) is usually traced to Persian \emph{darkhwāst}, \textquotedblleft{}a request\textquotedblright{}.
\end{cousinweb}

\begin{grammarlens}[title={What you've built}]
One more word for giving reasons and opinions.
\end{grammarlens}

\subsection*{Your turn}
\begin{itemize}
\raggedright
  \item \emph{Say it:} \emph{darkhāst}
  \item \emph{Say it:} \emph{darkhāst}, once more
  \item \emph{Say it:} \emph{maiṅ ne naukrī ke liye darkhāst dī.} — I applied for a job.
  \item \emph{Recall:} say the Urdu for to impose, then the Urdu for a government department, then say \emph{darkhāst} again
\end{itemize}

\begin{tcolorbox}[breakable,colback=teal!4,colframe=teal!35!black,title={Before you move on}]
What does \ur{درخواست} mean? (\textquotedblleft{}An application; a request\textquotedblright{}.) Say it once more.
\end{tcolorbox}
\section[\texorpdfstring{\ur{عہدہ} (uhdā)}{عہدہ (uhdā)}]{\ur{عہدہ} (uhdā) — a post; a rank (in an office)}

\label{lesson:UR-C343-uhda}



\begin{quote}
Before the new one: say the Urdu for a government department, then the Urdu for an application.
\end{quote}

\subsection*{You'll want to know: \ur{عہدہ}}
\textbf{\ur{عہدہ}} — \emph{uhdā} — \textquotedblleft{}a post; a rank (in an office)\textquotedblright{}.

\begin{quote}
\ur{وہ} \ur{ایک} \ur{بڑے} \ur{عہدے} \ur{پر} \ur{ہے۔}
\end{quote}

\begin{quote}
\emph{voh ek baṛe uhde par hai.}
\end{quote}

\begin{quote}
She holds a senior post.
\end{quote}

Masculine.

\begin{cousinweb}
\textbf{\ur{عہدہ}} (\emph{uhdā}) is usually traced, through Persian, to Arabic \emph{ʿuhda}, \textquotedblleft{}a charge, a responsibility\textquotedblright{}: the duty a post carries.
\end{cousinweb}

\begin{grammarlens}[title={What you've built}]
One more word for giving reasons and opinions.
\end{grammarlens}

\subsection*{Your turn}
\begin{itemize}
\raggedright
  \item \emph{Say it:} \emph{uhdā}
  \item \emph{Say it:} \emph{uhdā}, once more
  \item \emph{Say it:} \emph{voh ek baṛe uhde par hai.} — She holds a senior post.
  \item \emph{Recall:} say the Urdu for a government department, then the Urdu for an application, then say \emph{uhdā} again
\end{itemize}

\begin{tcolorbox}[breakable,colback=teal!4,colframe=teal!35!black,title={Before you move on}]
What does \ur{عہدہ} mean? (\textquotedblleft{}A post; a rank (in an office)\textquotedblright{}.) Say it once more.
\end{tcolorbox}
\section[\texorpdfstring{\ur{ہدایت} (hidāyat)}{ہدایت (hidāyat)}]{\ur{ہدایت} (hidāyat) — an instruction; guidance}

\label{lesson:UR-C343-hidayat}



\begin{quote}
Before the new one: say the Urdu for an application, then the Urdu for a post in an office.
\end{quote}

\subsection*{You'll want to know: \ur{ہدایت}}
\textbf{\ur{ہدایت}} — \emph{hidāyat} — \textquotedblleft{}an instruction; guidance\textquotedblright{}.

\begin{quote}
\ur{افسر} \ur{نے} \ur{ہمیں} \ur{ہدایت} \ur{دی۔}
\end{quote}

\begin{quote}
\emph{afsar ne hameṅ hidāyat dī.}
\end{quote}

\begin{quote}
The officer gave us instructions.
\end{quote}

Feminine.

\begin{cousinweb}
\textbf{\ur{ہدایت}} (\emph{hidāyat}) is usually traced, through Persian, to Arabic \emph{hidāya}, \textquotedblleft{}guidance\textquotedblright{}.
\end{cousinweb}

\begin{grammarlens}[title={What you've built}]
One more word for giving reasons and opinions.
\end{grammarlens}

\subsection*{Your turn}
\begin{itemize}
\raggedright
  \item \emph{Say it:} \emph{hidāyat}
  \item \emph{Say it:} \emph{hidāyat}, once more
  \item \emph{Say it:} \emph{afsar ne hameṅ hidāyat dī.} — The officer gave us instructions.
  \item \emph{Recall:} say the Urdu for an application, then the Urdu for a post in an office, then say \emph{hidāyat} again
\end{itemize}

\begin{tcolorbox}[breakable,colback=teal!4,colframe=teal!35!black,title={Before you move on}]
What does \ur{ہدایت} mean? (\textquotedblleft{}An instruction; guidance\textquotedblright{}.) Say it once more.
\end{tcolorbox}
\section[\texorpdfstring{\ur{سجھانا} (sujhānā)}{سجھانا (sujhānā)}]{\ur{سجھانا} (sujhānā) — to suggest; to put (an idea) to someone}

\label{lesson:UR-C343-sujhana}



\begin{quote}
Before the new one: say the Urdu for a post in an office, then the Urdu for an instruction.
\end{quote}

\subsection*{You'll want to know: \ur{سجھانا}}
\textbf{\ur{سجھانا}} — \emph{sujhānā} — \textquotedblleft{}to suggest; to put (an idea) to someone\textquotedblright{}.

\begin{quote}
\ur{اس} \ur{نے} \ur{مجھے} \ur{ایک} \ur{اچھا} \ur{راستہ} \ur{سجھایا۔}
\end{quote}

\begin{quote}
\emph{us ne mujhe ek achchhā rāstā sujhāyā.}
\end{quote}

\begin{quote}
She suggested a good way to me.
\end{quote}

Past, with \ur{نے} (he or she …): \ur{اس} \ur{نے} \ur{راستہ} \ur{سجھایا}, \emph{us ne rāstā sujhāyā}.

\begin{cousinweb}
\textbf{\ur{سجھانا}} (\emph{sujhānā}) is inherited: CDIAL lists the Hindi cousin \emph{sujhānā}, \textquotedblleft{}to explain\textquotedblright{}, under Sanskrit \emph{śúdhyati}, \textquotedblleft{}becomes clear (of doubts)\textquotedblright{} (CDIAL 12527). To make an idea clear to someone.
\end{cousinweb}

\begin{grammarlens}[title={What you've built}]
One more word for giving reasons and opinions.
\end{grammarlens}

\subsection*{Your turn}
\begin{itemize}
\raggedright
  \item \emph{Say it:} \emph{sujhānā}
  \item \emph{Say it:} \emph{sujhānā}, once more
  \item \emph{Say it:} \emph{us ne mujhe ek achchhā rāstā sujhāyā.} — She suggested a good way to me.
  \item \emph{Recall:} say the Urdu for a post in an office, then the Urdu for an instruction, then say \emph{sujhānā} again
\end{itemize}

\begin{tcolorbox}[breakable,colback=teal!4,colframe=teal!35!black,title={Before you move on}]
What does \ur{سجھانا} mean? (\textquotedblleft{}To suggest; to put (an idea) to someone\textquotedblright{}.) Say it once more.
\end{tcolorbox}
